\begin{proof}[Proof of \cref{obj:lem:published-fuzzy-testing}]
Fix a measurable decision \(\widehat\Psi\). All loss integrals and their suprema below take values in \([0,\infty]\).

\begin{enumerate}
\item Define the testing submodel, its component risks, and its worst-case risk by
\[
\begin{aligned}
\mathcal M_{\mathrm{test}}
&:=\{\mathbb B_z:z\text{ is an index}\}
  \cup\{\mathbb C_z:z\text{ is an index}\},\\
\mathcal R_{\mathrm{test}}(P)
&:=\int
  |\widehat\Psi(w)-\Psi(P)|\,dP^{\otimes v_0}(w),
  \qquad P\text{ a measure on the observation space},\\
\mathcal S_{\mathrm{test}}
&:=\sup_{P\in\mathcal M_{\mathrm{test}}}
  \mathcal R_{\mathrm{test}}(P).
\end{aligned}
\]
Every member of \(\mathcal M_{\mathrm{test}}\) is a probability measure. The finite-product maps
\(z\mapsto\mathbb B_z^{\otimes v_0}\) and
\(z\mapsto\mathbb C_z^{\otimes v_0}\) are probability kernels. Define their mixture laws by
\[
\mathbb B_{\mathrm{mix}}
:=\int\mathbb B_z^{\otimes v_0}\,d\omega(z),
\qquad
\mathbb C_{\mathrm{mix}}
:=\int\mathbb C_z^{\otimes v_0}\,d\omega(z).
\]
Both mixtures are probability measures, since their component laws and the prior are probability measures. The distance hypothesis therefore gives
\[
\mathsf h^2(\mathbb B_{\mathrm{mix}},\mathbb C_{\mathrm{mix}})
\le u_0.
\]
% lean: Causalean.Stat.Minimax.absolute_fuzzy_testing_lower_bound_iid

\item The probability prior ensures that the index space is nonempty. Choose an index \(z_0\). Separation gives
\[
\Psi(\mathbb C_z)\le \Psi(\mathbb B_{z_0})-s_0
\qquad\text{for every }z.
\]
Thus the nonempty set of values \(\Psi(\mathbb C_z)\) is bounded above. Define the finite real numbers
\[
\psi_{\mathbb C}^{\mathrm{sup}}
:=\sup_z\Psi(\mathbb C_z),
\qquad
a_{\mathrm{test}}
:=\psi_{\mathbb C}^{\mathrm{sup}}+\frac{s_0}{2}.
\]
For each index \(z\), separation with every \(z'\) gives
\[
\Psi(\mathbb C_{z'})
\le\Psi(\mathbb B_z)-s_0,
\qquad
\psi_{\mathbb C}^{\mathrm{sup}}
\le\Psi(\mathbb B_z)-s_0.
\]
Consequently,
\[
a_{\mathrm{test}}+\frac{s_0}{2}
\le\Psi(\mathbb B_z),
\qquad
\Psi(\mathbb C_z)
\le a_{\mathrm{test}}-\frac{s_0}{2}.
\]
Define the measurable threshold event on the sample space by
\[
\mathcal A_{\mathrm{test}}
:=\{w:\widehat\Psi(w)<a_{\mathrm{test}}\}.
\]
% lean: Causalean.Stat.Minimax.exists_separating_threshold

\item We derive the Hellinger testing comparison for these mixtures. Define
\[
\begin{aligned}
\xi_{\mathrm{test}}&:=\mathbb B_{\mathrm{mix}}+\mathbb C_{\mathrm{mix}},\\
f_{\mathbb B}(w)&:=\frac{d\mathbb B_{\mathrm{mix}}}{d\xi_{\mathrm{test}}}(w),
&
f_{\mathbb C}(w)&:=\frac{d\mathbb C_{\mathrm{mix}}}{d\xi_{\mathrm{test}}}(w),\\
\mathsf h_{\mathrm{test}}^2
&:=\mathsf h^2(\mathbb B_{\mathrm{mix}},\mathbb C_{\mathrm{mix}}).
\end{aligned}
\]
Choose nonnegative versions of the densities. They are integrable, and
\[
\int f_{\mathbb B}\,d\xi_{\mathrm{test}}
=\int f_{\mathbb C}\,d\xi_{\mathrm{test}}=1,
\qquad
\int(f_{\mathbb B}-f_{\mathbb C})\,d\xi_{\mathrm{test}}=0.
\]
For every measurable sample event \(\mathcal A\),
\[
\begin{aligned}
\mathbb B_{\mathrm{mix}}(\mathcal A)
-\mathbb C_{\mathrm{mix}}(\mathcal A)
&=\int_{\mathcal A}(f_{\mathbb B}-f_{\mathbb C})\,d\xi_{\mathrm{test}},\\
\int_{\mathcal A^c}(f_{\mathbb B}-f_{\mathbb C})\,d\xi_{\mathrm{test}}
&=-\int_{\mathcal A}(f_{\mathbb B}-f_{\mathbb C})\,d\xi_{\mathrm{test}}.
\end{aligned}
\]
Applying the integral triangle inequality on each of \(\mathcal A\) and
\(\mathcal A^c\), then adding, yields
\[
2\left|
\int_{\mathcal A}(f_{\mathbb B}-f_{\mathbb C})\,d\xi_{\mathrm{test}}
\right|
\le
\int|f_{\mathbb B}-f_{\mathbb C}|\,d\xi_{\mathrm{test}}.
\]

The normalization in \cref{obj:def:hellinger} gives
\[
\mathsf h_{\mathrm{test}}^2
=\int(\sqrt{f_{\mathbb B}}-\sqrt{f_{\mathbb C}})^2\,d\xi_{\mathrm{test}}.
\]
The square roots of the densities are square integrable. The pointwise identities
\[
\begin{aligned}
(\sqrt{f_{\mathbb B}}+\sqrt{f_{\mathbb C}})^2
+(\sqrt{f_{\mathbb B}}-\sqrt{f_{\mathbb C}})^2
&=2f_{\mathbb B}+2f_{\mathbb C},\\
|f_{\mathbb B}-f_{\mathbb C}|
&=|\sqrt{f_{\mathbb B}}-\sqrt{f_{\mathbb C}}|\,
  |\sqrt{f_{\mathbb B}}+\sqrt{f_{\mathbb C}}|
\end{aligned}
\]
therefore imply
\[
\int(\sqrt{f_{\mathbb B}}+\sqrt{f_{\mathbb C}})^2\,d\xi_{\mathrm{test}}
=4-\mathsf h_{\mathrm{test}}^2.
\]
Also, by nonnegativity and the pointwise bound
\((\sqrt{f_{\mathbb B}}-\sqrt{f_{\mathbb C}})^2
\le f_{\mathbb B}+f_{\mathbb C}\),
\[
0\le\mathsf h_{\mathrm{test}}^2
\le\int(f_{\mathbb B}+f_{\mathbb C})\,d\xi_{\mathrm{test}}
=2.
\]
Cauchy--Schwarz now gives
\[
\begin{aligned}
\int|f_{\mathbb B}-f_{\mathbb C}|\,d\xi_{\mathrm{test}}
&\le
\sqrt{\mathsf h_{\mathrm{test}}^2}\,
\sqrt{4-\mathsf h_{\mathrm{test}}^2}\\
&=\sqrt{\mathsf h_{\mathrm{test}}^2
                 (4-\mathsf h_{\mathrm{test}}^2)}.
\end{aligned}
\]
Combining the preceding bounds proves, for every measurable \(\mathcal A\),
\[
\begin{aligned}
|\mathbb B_{\mathrm{mix}}(\mathcal A)
-\mathbb C_{\mathrm{mix}}(\mathcal A)|
&\le\frac12
 \sqrt{\mathsf h_{\mathrm{test}}^2(4-\mathsf h_{\mathrm{test}}^2)}\\
&=\sqrt{\mathsf h_{\mathrm{test}}^2
                   (1-\mathsf h_{\mathrm{test}}^2/4)}.
\end{aligned}
\]
% lean: Causalean.Stat.Minimax.measure_event_gap_le_sharp_hellinger

\item Define the total variation distance of the mixtures by
\[
\mathsf V_{\mathrm{test}}
:=\sup_{\mathcal A\text{ measurable}}
|\mathbb B_{\mathrm{mix}}(\mathcal A)
-\mathbb C_{\mathrm{mix}}(\mathcal A)|.
\]
Taking the supremum in the preceding event bound gives
\[
\mathsf V_{\mathrm{test}}
\le
\sqrt{\mathsf h_{\mathrm{test}}^2
                   (1-\mathsf h_{\mathrm{test}}^2/4)}.
\]
The complementary-event identity implies
\[
\begin{aligned}
\mathbb B_{\mathrm{mix}}(\mathcal A_{\mathrm{test}})
+\mathbb C_{\mathrm{mix}}(\mathcal A_{\mathrm{test}}^c)
&=1+\mathbb B_{\mathrm{mix}}(\mathcal A_{\mathrm{test}})
       -\mathbb C_{\mathrm{mix}}(\mathcal A_{\mathrm{test}})\\
&\ge 1-\mathsf V_{\mathrm{test}}.
\end{aligned}
\]
Since \(0\le\mathsf h_{\mathrm{test}}^2\le u_0<2\),
\[
u_0(1-u_0/4)
-\mathsf h_{\mathrm{test}}^2(1-\mathsf h_{\mathrm{test}}^2/4)
=
\frac{(u_0-\mathsf h_{\mathrm{test}}^2)
      (4-u_0-\mathsf h_{\mathrm{test}}^2)}{4}
\ge0.
\]
Monotonicity of the square root consequently yields
\[
1-\sqrt{u_0(1-u_0/4)}
\le
\mathbb B_{\mathrm{mix}}(\mathcal A_{\mathrm{test}})
+\mathbb C_{\mathrm{mix}}(\mathcal A_{\mathrm{test}}^c).
\]
% lean: Causalean.Stat.Minimax.binary_test_error_ge_of_hellinger_budget

\item The threshold separation gives the pointwise loss bounds, for every index \(z\),
\[
\begin{aligned}
\frac{s_0}{2}\mathbf 1_{\mathcal A_{\mathrm{test}}}(w)
&\le|\widehat\Psi(w)-\Psi(\mathbb B_z)|,\\
\frac{s_0}{2}\mathbf 1_{\mathcal A_{\mathrm{test}}^c}(w)
&\le|\widehat\Psi(w)-\Psi(\mathbb C_z)|.
\end{aligned}
\]
Indeed, on \(\mathcal A_{\mathrm{test}}\) the decision lies below
\(a_{\mathrm{test}}\), while \(\Psi(\mathbb B_z)\) lies at least \(s_0/2\)
above it. On its complement the decision lies at or above the threshold,
while \(\Psi(\mathbb C_z)\) lies at least \(s_0/2\) below it. Outside the
respective events, the displayed left-hand sides vanish.

Integrating these inequalities under the component product laws gives
\[
\begin{aligned}
\frac{s_0}{2}\mathbb B_z^{\otimes v_0}(\mathcal A_{\mathrm{test}})
&\le\mathcal R_{\mathrm{test}}(\mathbb B_z)
\le\mathcal S_{\mathrm{test}},\\
\frac{s_0}{2}\mathbb C_z^{\otimes v_0}(\mathcal A_{\mathrm{test}}^c)
&\le\mathcal R_{\mathrm{test}}(\mathbb C_z)
\le\mathcal S_{\mathrm{test}}.
\end{aligned}
\]
Kernel event probabilities are measurable functions of the index. Integrating
their bounds against the probability prior therefore gives
\[
\begin{aligned}
\frac{s_0}{2}\mathbb B_{\mathrm{mix}}(\mathcal A_{\mathrm{test}})
&=\int\frac{s_0}{2}
  \mathbb B_z^{\otimes v_0}(\mathcal A_{\mathrm{test}})\,d\omega(z)
\le\int\mathcal S_{\mathrm{test}}\,d\omega
=\mathcal S_{\mathrm{test}},\\
\frac{s_0}{2}\mathbb C_{\mathrm{mix}}(\mathcal A_{\mathrm{test}}^c)
&=\int\frac{s_0}{2}
  \mathbb C_z^{\otimes v_0}(\mathcal A_{\mathrm{test}}^c)\,d\omega(z)
\le\int\mathcal S_{\mathrm{test}}\,d\omega
=\mathcal S_{\mathrm{test}}.
\end{aligned}
\]
Adding yields the extended nonnegative inequality
\[
\frac{s_0}{2}
\left\{
\mathbb B_{\mathrm{mix}}(\mathcal A_{\mathrm{test}})
+\mathbb C_{\mathrm{mix}}(\mathcal A_{\mathrm{test}}^c)
\right\}
\le2\mathcal S_{\mathrm{test}}.
\]
% lean: Causalean.Stat.Minimax.separated_mixture_errors_le_two_mul_sup

\item Combining the testing and risk inequalities gives
\[
\frac{s_0}{2}\left\{1-\sqrt{u_0(1-u_0/4)}\right\}
\le2\mathcal S_{\mathrm{test}}.
\]
The bracket is nonnegative under \(0\le u_0<2\). Using
\[
\frac{s_0}{2}\left\{1-\sqrt{u_0(1-u_0/4)}\right\}
=
2\left[
\frac{s_0}{4}\left\{1-\sqrt{u_0(1-u_0/4)}\right\}
\right]
\]
and cancelling the finite positive factor \(2\) in \([0,\infty]\), we obtain
\[
\frac{s_0}{4}\left\{1-\sqrt{u_0(1-u_0/4)}\right\}
\le\mathcal S_{\mathrm{test}}.
\]
Finally, every member of \(\mathcal M_{\mathrm{test}}\) equals
\(\mathbb B_z\) or \(\mathbb C_z\) for some index \(z\); either membership
hypothesis places that measure in \(\mathcal M\). Thus
\[
\mathcal M_{\mathrm{test}}\subseteq\mathcal M,
\qquad
\mathcal S_{\mathrm{test}}
\le
\sup_{P\in\mathcal M}
\int|\widehat\Psi(w)-\Psi(P)|\,dP^{\otimes v_0}(w).
\]
Transitivity proves the asserted bound, including when the supremum is infinite.
% lean: CausalSmith.Stat.TwosamplePointcateAnnotationFrontier.PublishedFuzzyTesting
\end{enumerate}
\end{proof}
